% =====================================================================
% CAPITULO 5 - GENERALIZACION Y ROBUSTEZ
% =====================================================================
% Responde la pregunta con la que cierra el capitulo 4: si dentro del dia las
% tres vistas saturan, hay que perturbar las condiciones y ver cual sigue en
% pie. Cuatro perturbaciones de exigencia creciente.
%
% ESTILO (31-ago): sin negritas en el texto corrido, sin incisos entre rayas
% salvo cuando la frase los pida, y sin florituras.
% =====================================================================

El capítulo anterior termina con una pregunta. Si dentro de un mismo día las
tres vistas aciertan por encima de $0.97$ y la comparación no las distingue,
cómo se decide cuál usar. La respuesta de este capítulo es que hay que perturbar
las condiciones y ver cuál sigue en pie.

Se aplican cuatro perturbaciones, de menor a mayor exigencia: mirar qué parte de
la entrada usa realmente el modelo, dejar que el atacante modifique esa parte,
aplicar el modelo a otro día del mismo régimen y aplicarlo a otro conjunto de
datos. Bajo todas ellas las vistas dejan de ser equivalentes, y el orden en que
caen no es el que sugerían sus métricas.

\section{Modelos empleados}
\label{sec:modelos}

Antes de perturbar queda declarar qué se perturba. Este trabajo emplea diez
modelos, y la distinción entre los propios y los tomados del estado del arte se
mantiene explícita en toda la memoria.

\subsection*{Modelos propios}

\begin{itemize}
  \item Perceptrón multicapa de dos capas ocultas, de 256 y 128 unidades, con
        activación ReLU, optimizador Adam, regularización $L_2$ de $10^{-4}$ y
        parada temprana, precedido de un escalado estándar ajustado solo sobre
        la partición de entrenamiento. Es el clasificador que produce todas las
        comparaciones entre vistas del capítulo~\ref{ch:regimenes}, y se emplea
        también sobre las representaciones tabulares del resto de la memoria.
        Su número de parámetros depende de la dimensión de entrada de cada
        vista, que va de 4 a 257, por lo que no entra en la comparación a
        presupuesto igualado que se describe más abajo.
  \item byte-CNN: red convolucional unidimensional sobre la secuencia cruda de
        bytes, con \emph{kernels} de varios tamaños en paralelo (3, 5 y 7) para
        captar $n$-gramas de longitud distinta, seguidos de agrupación máxima
        global. 29.762 parámetros.
  \item byte-LSTM: red recurrente bidireccional sobre la misma entrada, pensada
        como contraste, porque si la firma fuese secuencial de largo alcance
        debería ganar. 52.482 parámetros.
  \item Híbrido de dos ramas con fusión tardía: una rama convolucional sobre los
        bytes y otra densa sobre los metadatos, concatenadas en el espacio de
        representaciones y no en el de entradas. 31.042 parámetros.
  \item Autoencoder e Isolation Forest para la detección no supervisada,
        entrenados únicamente con tráfico benigno.
\end{itemize}

\subsection*{Modelos del estado del arte, implementados para contrastar}

\begin{itemize}
  \item byte-Transformer: codificador con autoatención sobre la secuencia de
        bytes, con \emph{token} \texttt{[CLS]} y codificación posicional
        aprendida, de 4 cabezas y 2 capas. 41.922 parámetros.
  \item CrossAttentionNet: fusión por atención cruzada bidireccional entre la
        rama de bytes y la de metadatos, en lugar de concatenación. 72.066
        parámetros.
  \item Random Forest y XGBoost \cite{Chen2016} sobre características de flujo,
        que son el estándar de facto en la literatura de detección de intrusiones
        tabular.
\end{itemize}

Los modelos clásicos se construyen sobre scikit-learn \cite{Pedregosa2011} y los
profundos sobre PyTorch \cite{Paszke2019}.

\subsection*{Por qué el presupuesto de parámetros importa}

Todos los modelos secuenciales se comparan en una banda estrecha, entre 29.762 y
72.066 parámetros. No es un detalle de implementación. Sin esa restricción, una
comparación entre arquitecturas mide sobre todo la capacidad: un
\emph{transformer} de diez millones de parámetros que superase a un byte-CNN de
treinta mil no demostraría que la autoatención sea mejor para este problema,
sino que un modelo mayor ajusta mejor, o peor si sobreajusta. Igualar el
presupuesto es lo que permite atribuir la diferencia a la arquitectura, que es
lo que se quiere medir.

\section{Interpretabilidad: qué mira el modelo}
\label{sec:saliency}

La primera perturbación no es tal. Consiste en abrir el modelo y ver qué parte
de la entrada usa, y es el paso que motiva todo lo demás.

Sobre el byte-CNN entrenado en el día del SSH-Bruteforce se calcula la saliencia
por gradientes: la magnitud del gradiente de la salida respecto a cada byte de la
entrada, promediada sobre los flujos de ataque. Mide cuánto cambiaría la decisión
si ese byte cambiase.

El resultado, en la figura~\ref{fig:saliency}, es inequívoco. El 31.1\% de la
importancia total se concentra en los primeros 48 bytes de un flujo que puede
llegar a 64 kilobytes. Esa zona es exactamente donde vive el \emph{banner} SSH
en claro, que precede a la negociación de la sesión. El cuerpo cifrado, que
ocupa el resto, apenas contribuye.

Inspeccionando esos bytes se ve por qué. El atacante emite dos banners distintos
en todos sus flujos, siendo el modal literalmente
\texttt{SSH-2.0-paramiko\_2.0.0}, frente a los 81 banners distintos de los
clientes legítimos, entre ellos OpenSSH, libssh2 y PuTTY. Y la entropía de la
carga útil es de $7.4$ en ambas clases, así que el contenido cifrado es
literalmente indistinguible.

El modelo no ve el cifrado, ve la herramienta. Es una firma real, porque el
atacante efectivamente usa \texttt{paramiko}, pero no es una firma del ataque
sino del programa con que se ejecutó. Cualquier atacante que cambie de
biblioteca la anula, sin modificar su técnica.

Este hallazgo es el que obliga a todo lo que sigue. Si un macro-F1 alto puede
sostenerse sobre una firma así, la métrica por sí sola no informa de nada, y hay
que someter cada vista a la pregunta de qué ocurre cuando esa firma desaparece.

\begin{figure}[htbp]
  \centering
  \includegraphics[width=\textwidth]{F24_saliency}
  \caption[Dónde mira el byte-CNN]{Importancia atribuida a cada posición de la
  entrada. El 31.1\% se concentra en los primeros 48 bytes, que es donde vive el
  \emph{banner} en claro; el cuerpo cifrado apenas contribuye.}
  \label{fig:saliency}
\end{figure}

\section{Robustez ante evasión}
\label{sec:evasion}

La forma directa de responder a esa pregunta es hacer lo que haría un atacante:
sobrescribir los primeros bytes de cada flujo de ataque con los de un cliente
legítimo, dejando intacto todo lo demás. Es una modificación barata, cambiar de
biblioteca o falsear una cabecera, y no altera la técnica del ataque.

\subsection*{Una advertencia metodológica que costó un experimento}

Un test de evasión solo es válido si el ataque sigue siendo el ataque después de
la modificación. Es fácil pasarlo por alto, y en este trabajo se pasó.

El primer experimento sobrescribía 160 bytes en todos los días. En el día web el
\emph{recall} del \emph{payload} se desplomaba de $0.9672$ a $0.0164$, lo
que parecía demostrar que también allí la firma era una huella de herramienta.
Inspeccionando los flujos modificados quedó claro que el resultado no
significaba eso: en 160 bytes cabe la inyección SQL entera, de modo que
sobrescribirlos no evade la detección, destruye el ataque. El modelo no fallaba,
acertaba al no marcar como ataque un flujo que ya no lo era.

Repetido con un prefijo de 24 bytes, que no llega a la carga maliciosa, el
\emph{payload} aguanta en $0.8852$. La conclusión se invierte por completo.

\subsection*{El resultado}

La tabla~\ref{tab:evasion-cap5} recoge el experimento corregido.

\begin{table}[htbp]
  \centering
  \caption[Robustez del payload ante camuflaje de los primeros bytes]{\emph{Recall}
  de la clase de ataque antes y después del camuflaje.}
  \label{tab:evasion-cap5}
  \begin{tabular}{lrr}
    \toprule
    \textbf{Día y ataque} & \textbf{Original} & \textbf{Camuflado} \\
    \midrule
    02-03 Bot (Ares)     & 1.0000 & 0.0025 \\
    20-02 DDoS-LOIC      & 1.0000 & 0.0108 \\
    16-02 DoS-Hulk       & 1.0000 & 0.0417 \\
    14-02 SSH-Bruteforce & 1.0000 & 0.3538 \\
    22-02 Web (24 B)     & 0.9672 & 0.8852 \\
    \bottomrule
  \end{tabular}
\end{table}

La tabla recoge solo el análisis de contenido, porque es el único al que el
camuflaje afecta: las vistas de metadatos y de conducta no leen un byte del
contenido y sus valores no cambian. La figura~\ref{fig:evasion} las incluye para
verlo.

Cuatro de los cinco días se desploman. En el día de la botnet el \emph{payload}
pasa de detectarlo todo a detectar el $0.25\%$, y en el DDoS al $1\%$. El
único que resiste es el día web, y resiste porque allí los bytes que el modelo
usa son el ataque: \texttt{union select} no es un rastro de la inyección SQL, es
la inyección.

Esto rescata un resultado que el capítulo anterior había dejado en suspenso. La
comparación de macro-F1 no distinguía el día web de los demás, porque todos
saturaban, pero la evasión sí lo distingue: es el único régimen en que el
análisis de contenido detecta el ataque y no la herramienta. Lo que la métrica
no separaba, la perturbación lo separa.

\subsection*{Invariante no es lo mismo que robusto}

Las vistas de metadatos y de conducta no cambian un solo valor bajo este ataque,
sencillamente porque no leen el contenido. Es lo que es:
una ventaja estructural frente a esta forma concreta de evasión, no una virtud
general. Un atacante que espacie deliberadamente sus conexiones ataca la ráfaga
con el mismo esfuerzo, y este trabajo no ha medido ese caso. Afirmar que la
conducta es robusta sin más sería el tipo de conclusión que el capítulo anterior
desmonta.

\begin{figure}[htbp]
  \centering
  \includegraphics[width=\textwidth]{F39_evasion_interarribo}
  \caption[El camuflaje separa la huella de herramienta del contenido
  real]{\emph{Recall} de la clase de ataque tras sobrescribir los primeros bytes
  de cada flujo, por vista y por día. El panel derecho compara los dos tamaños
  de prefijo ensayados en el día web.}
  \label{fig:evasion}
\end{figure}

\section{Cómo combinar las dos ramas}
\label{sec:fusion}

Antes de seguir perturbando queda resolver una cuestión de diseño que el
capítulo~\ref{ch:estado-del-arte} planteaba sin datos: si se dispone de una vista
de contenido y otra de flujo, dónde conviene unirlas.

Se comparan tres estrategias. La primera es la fusión ingenua, que concatena las
entradas crudas y entrega el vector resultante a un único modelo. La segunda es
la fusión tardía de dos ramas, en la que cada rama construye su representación
por separado y la concatenación ocurre en el espacio de representaciones, justo
antes de la capa de decisión. La tercera es cada rama por separado, como
referencia.

La figura~\ref{fig:arquitectura} muestra la segunda, que es la que este trabajo
emplea, y señala dónde habría concatenado la primera.

\begin{figure}[htbp]
  \centering
  \includegraphics[width=\textwidth]{F03_arquitectura}
  \caption[El híbrido de dos ramas]{El híbrido de dos ramas con fusión tardía.
  Cada vista se convierte primero en su propia representación y la unión ocurre
  entre representaciones de dimensión comparable, no entre entradas crudas.}
  \label{fig:arquitectura}
\end{figure}

\begin{table}[htbp]
  \centering
  \caption[Estrategias de fusión de las dos vistas]{Macro-F1 \emph{out-of-fold}
  de cada rama por separado y de las dos estrategias de fusión, en los dos
  regímenes opuestos y en la tarea multitipo del día web.}
  \label{tab:fusion}
  \footnotesize
  \begin{tabular}{lrrr}
    \toprule
    \textbf{Estrategia} & \textbf{14-02 SSH} & \textbf{22-02 web} & \textbf{22-02 multitipo} \\
    \midrule
    Solo payload (byte-CNN)   & 0.9988 & 1.0000 & 0.9679 \\
    Solo metadatos (MLP)      & 0.9931 & 0.9184 & 0.7695 \\
    Fusión ingenua (concatenar entradas) & 0.9988 & 0.9803 & 0.9009 \\
    Fusión tardía (dos ramas) & 1.0000 & 0.9975 & 0.9628 \\
    \bottomrule
  \end{tabular}
\end{table}

La tabla~\ref{tab:fusion} deja ver el problema de la fusión ingenua. En el día
web pasa de $1.0000$, que es lo que consigue la rama de contenido sola, a
$0.9803$; y en la tarea multitipo, de $0.9679$ a $0.9009$. Fusionar las
entradas hace el resultado peor que la mejor de sus partes.

El motivo es de escala. El histograma aporta 257 dimensiones, que son las 256
frecuencias de byte más la entropía, y los metadatos aportan cuatro escalares, de
modo que al concatenarlos los segundos quedan diluidos y el modelo
acaba comportándose como si solo tuviera la vista de contenido, con la
desventaja añadida de un espacio de entrada mayor y peor condicionado. La fusión
tardía evita el problema porque cada rama se entrena para producir una
representación de dimensión comparable antes de unirlas.

Es un detalle de diseño con una consecuencia práctica: la ventaja del enfoque
híbrido no se obtiene por el hecho de disponer de las dos fuentes, sino por
unirlas en el sitio correcto. Un trabajo que concatene entradas y concluya que la
fusión no aporta nada habrá medido la implementación, no la idea.

\section{Detección sin etiquetas de ataque}
\label{sec:anomalia}

Todo lo anterior supone un conjunto de entrenamiento con ataques etiquetados. Esa
suposición es fuerte, porque un ataque nuevo, por definición, no está etiquetado
en ninguna parte. La alternativa es la detección de anomalías: entrenar
únicamente con tráfico benigno y marcar lo que se desvía.

Se ensayan dos métodos clásicos, un Isolation Forest y un autoencoder, sobre las
mismas vistas y en tres regímenes. Se reporta el área bajo la curva ROC, porque
en este planteamiento no hay umbral natural y la métrica debe ser independiente
de él.

\begin{table}[htbp]
  \centering
  \caption[Detección de anomalías entrenada solo con tráfico benigno]{Área bajo
  la curva ROC. IF es Isolation Forest y AE es autoencoder. El azar corresponde a
  $0.5$.}
  \label{tab:anomalia}
  \begin{tabular}{lrrr}
    \toprule
    \textbf{Régimen} & \textbf{Payload (IF)} & \textbf{Payload (AE)} & \textbf{Conducta (IF)} \\
    \midrule
    Cifrado (14-02, SSH)     & 0.6946 & 0.5683 & 1.0000 \\
    En claro (22-02, web)    & 0.7580 & 0.7920 & 0.9121 \\
    Volumétrico (16-02, DoS) & 0.3411 & 0.4791 & 0.5402 \\
    \bottomrule
  \end{tabular}
\end{table}

La tabla~\ref{tab:anomalia} dice tres cosas. La primera es que sin etiquetas de
ataque la tesis del trabajo se replica: cada régimen esconde su anomalía en una
vista distinta, y la conducta resuelve el cifrado con $1.0000$ mientras el
contenido se queda en $0.6946$.

La segunda es que el régimen volumétrico es intrínsecamente difícil sin
etiquetas. Las tres combinaciones quedan por debajo de $0.55$, y el
Isolation Forest sobre contenido baja de $0.5$, es decir, ordena peor que el
azar. El motivo es que las peticiones que genera Hulk son sintácticamente más
regulares que el tráfico web legítimo, y un detector de anomalías
entrenado con lo normal las considera menos anómalas que el propio tráfico
normal.

La tercera es una advertencia sobre las cifras del resto del capítulo. Los
valores supervisados que aparecen en toda la memoria por encima de $0.97$
dependen de disponer de ataques etiquetados del mismo tipo. Cuando esa
suposición se retira, el rendimiento cae entre veinte y cincuenta puntos según el
régimen. Es la distancia real entre lo que se mide en un conjunto de datos
etiquetado y lo que un detector encontraría ante una amenaza nueva.

\section{Generalización entre días del mismo régimen}
\label{sec:cross-dia}

La segunda perturbación es cambiar de día. Se toman los tres días volumétricos
del 2018, que son el 16-02, el 15-02 y el 20-02, y se forman los seis pares
dirigidos posibles: se entrena en uno y se evalúa en otro, con el mismo régimen
de ataque pero distinta herramienta, distinta víctima y distinto tráfico de
fondo.

\subsection*{Una lección metodológica antes que un resultado}

Este hallazgo se enunció tres veces y cambió las tres. Merece contarse así
porque la lección vale tanto como el número.

La primera versión, con dos días, afirmaba que el modelo más expresivo
generaliza peor, porque el byte-CNN colapsaba y el histograma no. La segunda, al
añadir el tercer día y completar los seis pares, mostró que la regla era falsa:
se había inducido una ley con dos observaciones. La tercera llegó al cambiar de
métrica, y mostró que lo que colapsaba no era lo que parecía.

\subsection*{El diagnóstico: no falla la representación, falla el umbral}

Con umbral fijo en $0.5$, el byte-CNN da un macro-F1 de $0.3333$ en cuatro
de los seis pares. Ese valor no es simplemente malo: es exactamente lo que
produce un clasificador que asigna todo a una sola clase. Su \emph{recall} de
ataque es $0.0000$. Leído así, el modelo no transfiere en absoluto.

El AUC dice otra cosa. La tabla~\ref{tab:calibracion} lo pone al lado.

\begin{table}[htbp]
  \centering
  \caption[Transferencia entre días: umbral fijo frente a AUC]{Entrenar en un
  día volumétrico y evaluar en otro. El macro-F1 con umbral fijo sugiere
  colapso; el AUC muestra que el orden de las puntuaciones se conserva.}
  % @fuente [x for p in ["16-02 -> 20-02", "20-02 -> 16-02", "15-02 -> 20-02", "20-02 -> 15-02", "15-02 -> 16-02", "16-02 -> 15-02"] for x in R.CALIBRACION_CRUZADA[p]["byte-CNN"]]
  \label{tab:calibracion}
  \begin{tabular}{lrrr}
    \toprule
    \textbf{Par (entreno $\rightarrow$ prueba)} & \textbf{AUC-ROC} & \textbf{F1 (umbral 0.5)} & \textbf{F1 (umbral óptimo)} \\
    \midrule
    16-02 $\rightarrow$ 20-02 & 0.9724 & 0.0000 & 0.9849 \\
    20-02 $\rightarrow$ 16-02 & 0.9886 & 0.0000 & 0.9918 \\
    15-02 $\rightarrow$ 20-02 & 0.9758 & 0.0000 & 0.9872 \\
    20-02 $\rightarrow$ 15-02 & 0.9854 & 0.0000 & 0.9899 \\
    15-02 $\rightarrow$ 16-02 & 1.0000 & 1.0000 & 1.0000 \\
    16-02 $\rightarrow$ 15-02 & 1.0000 & 0.9581 & 0.9999 \\
    \bottomrule
  \end{tabular}
\end{table}

La figura~\ref{fig:calibracion} contrapone ambas lecturas. El AUC no baja de
$0.9724$ en ninguno de los seis pares. El modelo ordena correctamente los
flujos del día nuevo, porque los de ataque reciben puntuaciones más altas que
los benignos, y sin embargo, con el umbral heredado del día de entrenamiento, no
acierta ni uno. La representación transfiere; lo que no transfiere es dónde
cortar.

Es una distinción con consecuencias prácticas inmediatas, y es invisible para
cualquier métrica que dependa de un umbral fijo. Un trabajo que hubiese
reportado solo el F1 habría concluido que el byte-CNN no generaliza entre días,
y habría estado equivocado.

\subsection{De diagnóstico a receta: recalibración}

Si el problema es el umbral, la solución es recalibrar. Se aplica escalado de
Platt \cite{Platt1999}, que ajusta una función
$\sigma(A \cdot \mathrm{logit}(p) + B)$ sobre unas pocas muestras etiquetadas
del día de destino, sin reentrenar el modelo. La pregunta operativa es cuántas
etiquetas hacen falta, y la respuesta son veinticinco. La
tabla~\ref{tab:recalibracion} recoge los cuatro pares que colapsaban.

\begin{table}[htbp]
  \centering
  \caption[Recalibración con 25 etiquetas del día de destino]{Macro-F1 y error
  de calibración esperado (ECE) antes y después de aplicar escalado de Platt con
  25 flujos etiquetados del día de destino. No se reentrena el modelo.}
  % @fuente [x for p in ["16-02 -> 20-02", "20-02 -> 16-02", "15-02 -> 20-02", "20-02 -> 15-02"] for x in (R.RECALIBRACION[p]["sin"][0], R.RECALIBRACION[p]["sin"][2], R.RECALIBRACION[p]["n25"][0], R.RECALIBRACION[p]["n25"][2])]
  \label{tab:recalibracion}
  \begin{tabular}{lrrrr}
    \toprule
    & \multicolumn{2}{c}{\textbf{Sin recalibrar}} & \multicolumn{2}{c}{\textbf{Con 25 etiquetas}} \\
    \cmidrule(lr){2-3} \cmidrule(lr){4-5}
    \textbf{Par} & \textbf{F1} & \textbf{ECE} & \textbf{F1} & \textbf{ECE} \\
    \midrule
    16-02 $\rightarrow$ 20-02 & 0.3333 & 0.4973 & 0.9750 & 0.0216 \\
    20-02 $\rightarrow$ 16-02 & 0.3333 & 0.4991 & 0.9912 & 0.0087 \\
    15-02 $\rightarrow$ 20-02 & 0.3332 & 0.4967 & 0.9828 & 0.0170 \\
    20-02 $\rightarrow$ 15-02 & 0.3333 & 0.4993 & 0.9698 & 0.0216 \\
    \bottomrule
  \end{tabular}
\end{table}

La figura~\ref{fig:recalibracion} contrapone las dos situaciones en los seis
pares. El error de calibración cae de aproximadamente $0.499$,
que es el máximo posible y corresponde a un modelo sistemáticamente seguro y
sistemáticamente equivocado, a menos de $0.022$. Ampliar de 25 a 1.000
etiquetas apenas mejora nada, salvo en el par más difícil, donde el
\emph{recall} sube de $0.9571$ a $0.9997$.

Traducido a despliegue: cuando un detector entrenado en unas condiciones se
lleva a otras, no hay que reentrenarlo, basta con que un analista etiquete
veinticinco conexiones. Es una recomendación concreta, barata y verificable, y
sale directamente del diagnóstico anterior.

El mismo fenómeno ha sido descrito de forma independiente en la literatura
reciente \cite{CrossDatasetTransformer2026}, que lo corrige también con escalado
de Platt. Que dos trabajos lleguen por separado al mismo diagnóstico refuerza
que no es un artefacto de este montaje.

\begin{figure}[htbp]
  \centering
  \includegraphics[width=\textwidth]{F38_calibracion_cruzada}
  \caption[Lo que el AUC revela y el umbral fijo esconde]{Para cada par de días,
  el macro-F1 con umbral fijo frente al AUC. El primero sugiere colapso; el
  segundo muestra que el orden de las puntuaciones se conserva.}
  \label{fig:calibracion}
\end{figure}

\begin{figure}[htbp]
  \centering
  \includegraphics[width=\textwidth]{F41_recalibracion}
  \caption[Recalibración con pocas etiquetas]{Macro-F1 y error de calibración
  antes y después de recalibrar con 25 flujos etiquetados del día de destino, en
  los seis pares de días. Los cuatro pares que colapsaban se recuperan y el error
  de calibración cae de su valor máximo a casi cero.}
  \label{fig:recalibracion}
\end{figure}

\section{Una vista para la periodicidad, y qué la hace funcionar}
\label{sec:interarribo}

El capítulo~\ref{ch:regimenes} deja una promesa abierta. La firma de la botnet
vive en el intervalo entre conexiones, y ninguna de las tres vistas del protocolo
lo calcula: la ráfaga cuenta cuántas conexiones hay en una ventana, no cada
cuánto llegan. Esta sección cierra esa promesa añadiendo una cuarta vista y
midiendo qué aporta.

La vista de inter-arribo son seis características por flujo, todas derivadas de
los intervalos entre conexiones consecutivas del mismo origen: el intervalo
anterior y el siguiente, la media y el coeficiente de variación en una ventana
local, una medida de regularidad y la fracción de intervalos cercanos a la
mediana. Se calculan desde el mismo fichero que las demás, sin volver a procesar
las capturas.

\subsection*{Qué aporta como grupo}

\begin{table}[htbp]
  \centering
  \caption[Aportación de cada grupo de características]{Macro-F1
  \emph{out-of-fold} en el día de la botnet, combinando los metadatos de flujo
  con cada grupo temporal y con los grupos por separado.}
  \label{tab:ablacion-grupos}
  \begin{tabular}{lr}
    \toprule
    \textbf{Conjunto de características} & \textbf{Macro-F1} \\
    \midrule
    Metadatos de flujo                   & 0.9985 \\
    Metadatos + ráfaga                   & 0.9991 \\
    Metadatos + inter-arribo             & 0.9989 \\
    Metadatos + ráfaga + inter-arribo    & 0.9989 \\
    \midrule
    Inter-arribo solo (6 características) & 0.9914 \\
    Ráfaga sola (3 características)       & 0.9909 \\
    \bottomrule
  \end{tabular}
\end{table}

La tabla~\ref{tab:ablacion-grupos} obliga a una conclusión modesta y hay que
darla como tal. Sobre los metadatos, la ráfaga y el inter-arribo aportan
prácticamente lo mismo, con $0.9991$ frente a $0.9989$, y combinarlos no
mejora. Es decir, la vista nueva no supera a la que ya había.

Su interés está en otro sitio. Seis características temporales por sí solas, sin
mirar un byte de contenido ni un escalar de tamaño, alcanzan $0.9914$ en el día
de la botnet. Eso confirma que la periodicidad es una firma real y suficiente, y
que el trabajo la estaba detectando por otra vía. Y como no leen el contenido,
son invariantes ante el camuflaje de bytes de la sección~\ref{sec:evasion}, igual
que la ráfaga.

\subsection*{Qué característica hace el trabajo, y una corrección}

La ablación individual, en la tabla~\ref{tab:ablacion-inter}, mide dos cosas por
cada característica: cuánto se pierde al quitarla del conjunto y cuánto consigue
ella sola.

\begin{table}[htbp]
  \centering
  \caption[Ablación de las características de inter-arribo]{Para cada
  característica, macro-F1 del conjunto sin ella, caída resultante, y macro-F1
  que consigue por sí sola. Las seis juntas dan $0.9914$.}
  \label{tab:ablacion-inter}
  \begin{tabular}{lrrr}
    \toprule
    \textbf{Característica} & \textbf{Sin ella} & \textbf{Caída} & \textbf{Ella sola} \\
    \midrule
    \texttt{dt\_prev}            & 0.9896 & $+0.0018$ & 0.9669 \\
    \texttt{dt\_next}            & 0.9900 & $+0.0014$ & 0.9641 \\
    \texttt{media\_local}        & 0.9789 & $+0.0125$ & 0.9638 \\
    \texttt{cv\_local}           & 0.9835 & $+0.0079$ & 0.9355 \\
    \texttt{regularidad}         & 0.9949 & $-0.0035$ & 0.9062 \\
    \texttt{frac\_cerca\_mediana} & 0.9887 & $+0.0026$ & 0.9336 \\
    \bottomrule
  \end{tabular}
\end{table}

El resultado corrigió una afirmación previa de este trabajo. Se había escrito que
el coeficiente de variación local era la característica del \emph{beaconing},
porque mide precisamente la regularidad de los intervalos y esa es la intuición
evidente. La ablación lo desmiente por partida doble: la característica cuya
ausencia más cuesta es la media local, con una caída de $0.0125$ frente a
$0.0079$, y la que mejor funciona en solitario es el intervalo anterior en
crudo, con $0.9669$.

Hay además una característica que estorba. Quitar \texttt{regularidad} mejora el
resultado en $0.0035$, y en solitario es la peor de las seis con $0.9062$.
Es redundante con el coeficiente de variación y añade ruido, así que debería
eliminarse del conjunto.

La lectura de conjunto es que la señal está en los intervalos mismos y no en los
estadísticos que se construyen sobre ellos. Un solo intervalo en crudo alcanza
$0.9669$, y las cinco características derivadas añaden en total dos puntos y
medio. Es un recordatorio útil frente a la tentación de multiplicar
características derivadas: cada una necesita medirse antes de
incorporarla, porque alguna puede restar.

\section{La fuga de identificadores de red}
\label{sec:fuga}

La tercera perturbación descubre un problema que las dos anteriores no podían
ver, y que además sobrevive a la validación entre días.

Si a la vista de metadatos se le añaden las direcciones IP y el puerto de
origen, algo que buena parte de la literatura hace porque están disponibles en el
registro de flujo, el resultado entre días se vuelve perfecto: macro-F1 de
$1.0000$ en cinco de los seis pares con Random Forest y con XGBoost.

Parece una mejora y es un atajo, como cuantifica la figura~\ref{fig:fuga}. El
mecanismo se ve al mirar el direccionamiento del laboratorio: de las 18
direcciones atacantes que aparecen en los seis días, 15 empiezan por
\texttt{18.}, y las 12 víctimas están todas en \texttt{172.31.69.x}. Un modelo
que aprenda que un origen \texttt{18.x} implica ataque transfiere perfectamente
entre cualquier par de días del conjunto, sin haber aprendido nada sobre
ataques.

La parte que importa metodológicamente es que la validación entre días no
detecta el atajo. Al contrario, lo premia. Solo se rompe al cambiar de conjunto
de datos, donde el direccionamiento es distinto, y allí la vista de metadatos
cae al azar. Es el mismo problema que Engelen y sus coautores documentan en
CIC-IDS2017 \cite{Engelen2021}, aquí medido sobre CSE-CIC-IDS2018 y con una
consecuencia añadida: ninguna de las validaciones habituales dentro de un mismo
conjunto de datos lo detecta.

Un detalle refuerza el riesgo. Dos direcciones se reutilizan entre días con
papeles distintos: \texttt{18.219.211.138} es atacante de DoS-GoldenEye el 15-02
y canal de mando de la botnet el 02-03, y \texttt{18.218.115.60} ejecuta los
ataques web el 22-02 y es uno de los diez orígenes del DDoS el 20-02. Un modelo
que memorice direcciones acertará en ambos días por el motivo equivocado.

Por eso este trabajo excluye las direcciones y los puertos de todas las vistas
del capítulo~\ref{ch:regimenes}, y por eso ese mismo capítulo incluye la
dirección IP únicamente como control positivo declarado, nunca como
característica.

\section{Generalización entre conjuntos de datos}
\label{sec:cross-dataset}

La cuarta perturbación es la más exigente: entrenar en CSE-CIC-IDS2018 y evaluar
en CIC-IDS2017, sobre el ataque análogo. Cambia la red, el direccionamiento, el
tráfico de fondo y el año. Por indicación de la dirección del trabajo esta línea
quedó como validación complementaria, pero el capítulo la necesita, porque es la
única que rompe la fuga de la sección anterior.

Entrenando en 2018 y evaluando en 2017, el byte-CNN alcanza $0.9914$ en el
régimen web, $0.9990$ en el cifrado y $0.9656$ en el volumétrico, mientras
la vista de metadatos cae a $0.5287$, $0.5139$ y $0.4291$, es decir, al
azar. La lectura inmediata sería que el contenido generaliza y los metadatos no.

Después de la sección~\ref{sec:saliency} esa lectura ya no basta, y hay que
preguntar por qué transfiere cada caso:

\begin{itemize}
  \item El régimen web transfiere por contenido intrínseco. Las cadenas
        \texttt{union select} y \texttt{<script>} aparecen en ambos conjuntos
        porque son el ataque. Es transferencia robusta.
  \item El cifrado transfiere porque los atacantes de ambos años usan la misma
        biblioteca y emiten el mismo banner \texttt{paramiko}. Es un
        \emph{fingerprint} compartido, y por tanto evadible.
  \item El volumétrico, igual, por los \emph{User-Agent} y las cabeceras de Hulk
        y GoldenEye, comunes a ambos años.
\end{itemize}

Dos de los tres transfieren por el motivo equivocado. Un \emph{fingerprint}
compartido entre conjuntos de datos sigue siendo un \emph{fingerprint}: que
generalice no lo hace robusto, solo lo hace consistentemente frágil.

\subsection*{Una asimetría honesta}

La dirección contraria no funciona. Entrenando en 2017 y evaluando en 2018, el
byte-CNN colapsa: $0.5000$ en el régimen web, donde predice una sola clase,
$0.6051$ en el cifrado y $0.6710$ en el volumétrico. El motivo es la
diversidad, porque los ataques de 2018 son más numerosos y más variados, de modo
que un modelo entrenado allí ha visto más y transfiere, mientras que al revés no.

Es un argumento a favor de entrenar con la mayor variedad disponible, y también
un recordatorio de que la generalización cruzada no es simétrica: reportar solo
la dirección que funciona sería una forma elegante de mentir.

\begin{figure}[htbp]
  \centering
  \includegraphics[width=\textwidth]{F40_fuga_ips}
  \caption[El tamaño del atajo de las direcciones IP]{Macro-F1 entre días al
  añadir las direcciones y el puerto de origen a la vista de metadatos. La
  validación entre días no detecta el atajo: lo premia.}
  \label{fig:fuga}
\end{figure}

\section{¿Mejoran las arquitecturas del estado del arte?}
\label{sec:sota-contraste}

Queda una pregunta que atraviesa todo el trabajo: si los resultados no son
mejores, cabe preguntarse si es porque los modelos empleados se quedan cortos.
Para responderla se implementaron dos arquitecturas del estado del arte y se
compararon con las propias a presupuesto de parámetros comparable.

\subsection*{El transformer nunca gana}

Sobre la secuencia de bytes se comparan byte-CNN, byte-LSTM y byte-Transformer.
El resultado depende de la cantidad de datos y no favorece al modelo más
moderno. Con pocos datos, en el día web con 203 flujos por clase, el
\emph{transformer} es peor: $0.9778$ frente a $0.9975$ del byte-CNN. Con
datos suficientes, en el día de la botnet con 4.000 por clase, empata en
$1.0000$. Y siempre cuesta entre 36 y 49 veces más tiempo, con 3.191 segundos
frente a 65 en el día de la botnet.

La figura~\ref{fig:arquitecturas} contrapone rendimiento y coste. El byte-LSTM,
por su parte, se hunde con pocos datos hasta $0.8296$, lo que es coherente con
lo que mostró la interpretabilidad: la firma vive en los primeros bytes, y una
convolución con agrupación máxima la localiza mejor que una recurrencia que ha
de recorrer toda la secuencia.

\subsection*{La atención cruzada empata usando el doble de parámetros}

En la fusión de las dos ramas, la atención cruzada iguala a la fusión tardía,
con $0.9828$ frente a $0.9828$ en el día web y $0.9999$ frente a
$1.0000$ en el de la botnet, empleando 72.066 parámetros frente a 31.042, es
decir, 2.3 veces más. No hay ganancia que justifique el coste.

\subsection*{Y no es por falta de ajuste}

Cabría objetar que los modelos propios están mejor ajustados que los tomados de
la literatura. Se comprobó con una búsqueda aleatoria de hiperparámetros de 30
configuraciones, validada únicamente dentro del día de entrenamiento para no
contaminar la evaluación. La carencia era real pero pequeña: el mejor ajuste
mejora el MLP en $+0.0275$ en el día web y en $+0.0063$ en el volumétrico, y
a XGBoost en $+0.0122$ y $+0.0005$. Nada de eso cambia ninguna conclusión.

Esa comprobación tiene un alcance acotado. La búsqueda cubre los dos
modelos tabulares, el perceptrón multicapa y XGBoost, sobre los que se apoya la
mayor parte de las mediciones del trabajo. Las cinco redes profundas conservan
hiperparámetros fijos: el argumento vale para el ajuste de
los modelos tabulares y no es una búsqueda exhaustiva sobre las arquitecturas
profundas. Lo que sí queda igualado entre ellas, que es lo que la comparación
del apartado mide, es el presupuesto de parámetros.

\subsection*{Qué se sigue de esto}

Ninguna de las arquitecturas avanzadas mejora los resultados de este trabajo
sobre estos datos, y dos de ellas cuestan entre 2 y 49 veces más. No es un
argumento contra el estado del arte, cuyos modelos están pensados para escalas y
escenarios mayores, y esta comparación no los refuta en general.

Es, en cambio, la confirmación de lo que el capítulo~\ref{ch:regimenes} había
medido por otra vía. Si dentro del día la evaluación está saturada, no queda
margen que una arquitectura mejor pueda recoger. El cuello de botella de este
problema no está en la capacidad del modelo, sino en la calidad de la etiqueta,
en los atajos que el conjunto de datos permite y en cómo se mide la
transferencia. Que es donde este trabajo ha intervenido.

\begin{figure}[htbp]
  \centering
  \includegraphics[width=\textwidth]{F42_arquitecturas}
  \caption[Arquitecturas propias frente a las del estado del arte]{Macro-F1 y
  coste computacional a presupuesto de parámetros comparable. El
  \emph{transformer} nunca gana y siempre cuesta más.}
  \label{fig:arquitecturas}
\end{figure}
